\section{A boundary deformation in four dimensions}
\label{four:sec:output}

The construction in this section keeps a genuine inner boundary throughout
the graph deformation. The boundary has two jobs: it must become strictly
mean concave, so that a minimal enclosure can detach from it, and its flux
must be small enough to leave the energy at infinity almost unchanged.
These requirements explain the signed collars and the boundary condition
in the coupled equations. They also distinguish this construction from a
solve on a manifold obtained by filling the boundary.

We first state the geometric output and deduce its numerical consequence.
The rest of this part constructs that output. The proof is independent
of the numerical spacetime inequality: its scalar identity and local
analytic estimates apply before any Riemannian comparison is made.

\subsection{Data, normalization and the full enclosing area}

Here the spatial dimension is four. Write
\[
 \omega_3=|S^3|=2\pi^2,\qquad
 \tau=\tr_gK,\qquad
 \mu=\frac12(R_g+\tau^2-|K|_g^2),\qquad
 J_i=\nabla^j(K_{ij}-\tau g_{ij}).
\]
Here $\tau$ denotes the trace of $K$; the small height regularizer called
$\tau$ in the three-dimensional construction is denoted by $\eta_N$ in
this part. The tensor $K$ is symmetric, and $\nabla$ is the Levi--Civita
connection of $g$. Our Laplacian is $\Delta_g=\divg_g\grad_g$. A hypersurface with
unit normal $\nu$ has mean curvature $H=\divg_{\mathrm{tan}}\nu$ and
future expansion $H+\tr_{\mathrm{tan}}K$. At the inner boundary the
normal points into the exterior. In a compact auxiliary region it points
out of that region. Every change of sign below specifies both the tensor
and the orientation.

The notation $O_j(r^{-a})$ includes coordinate derivatives of order
$k\le j$, bounded by $O(r^{-a-k})$. In the Euclidean end coordinates our
energy and momentum conventions are
\begin{align}
 E_g&=\frac1{6\omega_3}\lim_{R\to\infty}
 \int_{|x|=R}(\partial_jg_{ij}-\partial_i g_{jj})
                \nu_\delta^i\,\dd A_\delta,
                   \label{four:intro:energy}\\
 P_i&=\frac1{3\omega_3}\lim_{R\to\infty}
 \int_{|x|=R}(K_{ij}-\tau g_{ij})\nu_\delta^j\,\dd A_\delta.
                   \label{four:intro:momentum}
\end{align}
The prepared tensor below has compact support, hence $P=0$. No additional
normalization of mass accompanies these density conventions.

\begin{definition}[Full cuts in a one-ended exterior]\label{four:def:cuts}
Let $M$ be a smooth connected manifold with nonempty compact boundary $S$
and one Euclidean end with compact complement. A full cut is the entire
intrinsic boundary of a smooth connected codimension-zero outer domain
$D\subset M$ that is closed in $M$, contains the entire distant end,
has compact boundary, and has its manifold interior in $\operatorname{int}M$.
Its boundary is embedded and two-sided; all components and every portion
coincident with $S$ count. Set
\[
 A_* = \inf_D |\partial_{\mathrm{man}}D|_g.
\]
In particular $D=M$ is permitted. The infimum need not be attained.
\end{definition}

These are full cuts rather than a selected connected component of a cut.
Filling bounded pockets on the obstacle side removes irrelevant internal
frontiers, while leaving a connected exterior containing the end. The
metric comparison in the construction will apply on every tangent
three-plane, so it will control this entire class at once.

\subsection{The geometric output and its immediate use}

\begin{theorem}[Four-dimensional boundary deformation]
\label{four:thm:boundary-deformation}
Let $(M^4,g,K)$ be smooth, connected and orientable, complete with its
nonempty compact smooth boundary $S$ included, with exactly one Euclidean
coordinate end and compact complement. Suppose $K$ is compactly supported.
Let $r$ be a smooth positive extension of the coordinate radius. Assume
that, for some $c_0>0$ and $0<\delta_1<1$,
\[
 \mu-|J|_g\ge c_0r^{-4-\delta_1}\quad\hbox{on }M,
 \qquad H+\tr_SK<0\quad\hbox{on every component of }S,
\]
with the normal into $M$. On the end assume
\[
 g-\delta=O_j(r^{-2})\quad\hbox{for every fixed }j\ge0,
 \qquad R_g=O(r^{-4-\delta_1}),
\]
and finite ADM energy in Equation~\eqref{four:intro:energy}.

There is a smooth closed connected one-ended exterior $\Omega\subset M$
with nonempty compact boundary $B\subset\operatorname{int}M$ such that
every full cut of $B$ in $\Omega$ has $g$-area at least $A_*$.
For every $0<\eps<1$ and all sufficiently large integers $N$ there is a
smooth metric $\widehat g_{\eps,N}$ on $\Omega$, complete with $B$
included, satisfying
\begin{gather}
 R_{\widehat g_{\eps,N}}\ge0,\qquad
 H_{\widehat g_{\eps,N}}(B)<0,
 \qquad \widehat g_{\eps,N}\ge e^{-2\eps}g,
                      \label{four:output:geometry}\\
 \widehat g_{\eps,N}-\delta=O_2(r^{-2}),\qquad
 R_{\widehat g_{\eps,N}}=O(r^{-4-\delta_1}),
                      \label{four:output:decay}\\
 E_{\widehat g_{\eps,N}}\le E_g+
 \frac{\Pi_N}{3\omega_3}
       \bigl(e^{-\eps N}+e^{-\eps N/2}\bigr).
                      \label{four:output:energy}
\end{gather}
The boundary mean curvature uses the normal into $\Omega$. The constant
$\Pi_N=C(1+N)^{a_{\mathrm{pol}}}$ is polynomial in $N$, with $C$ and
$a_{\mathrm{pol}}$ depending only
on the fixed data, $\eps$ and the fixed geometric choices. The constants
in Equation~\eqref{four:output:decay} may depend arbitrarily on $N$.
\end{theorem}

The metric has the form $e^{2t}(g+l(t)^2\dd f\otimes\dd f)$: the graph
term cannot shorten a tangent vector, and the floor $t\ge-\eps$ protects
all enclosing areas. Its inner boundary is not itself used in the
Riemannian inequality. Strictly negative mean curvature supplies an
outward barrier for a new minimal enclosure.

We use the following precise form of the independent smooth enclosure
theorem. It concerns perimeter alone and uses no stationary-field,
variation or equality conclusion.

\begin{proposition}[Smooth full-perimeter enclosure input]
\label{four:input:enclosure}
Let $(X^4,h)$ be a smooth connected orientable one-ended exterior, complete
with nonempty compact smooth boundary $B$ included. Suppose the end has
compact complement and $h-\delta=O_2(r^{-q})$ for some $q>1$. Minimize the
full perimeter of the filled obstacle among enclosing sets; contact with
$B$ and all frontier components count. The minimizing enclosure can be
chosen with connected complete exterior, and its full perimeter equals
the infimum over the smooth full cuts of $B$. Its exterior frontier is
outer area-minimizing against all enclosing competitors, including
disconnected ones. If $H_h(B)<0$ for the normal into $X$, that frontier
is a nonempty smooth compact embedded minimal hypersurface contained in
$\operatorname{int}X$.
\end{proposition}

This is the smooth enclosure theorem from the independent Riemannian
geometric preliminaries~\cite[Lemmas~2.1--2.3]{RiemannianCompanion}.
Its proof uses confinement, BV compactness,
smooth-cut approximation and obstacle regularity. Only its strict-boundary
case is used here.

\begin{corollary}[Energy bound for the prepared exterior]
\label{four:thm:prepared-energy}
Under the hypotheses of Theorem~\ref{four:thm:boundary-deformation},
\begin{equation}\label{four:ex:prepared-inequality}
 E_g\ge\frac12\left(\frac{A_*}{\omega_3}\right)^{2/3}.
\end{equation}
\end{corollary}

\begin{proof}
Fix $\eps>0$ and a sufficiently large $N$, and abbreviate
$\widehat g=\widehat g_{\eps,N}$. Proposition~\ref{four:input:enclosure}
applies to $(\Omega,\widehat g)$: Theorem~\ref{four:thm:boundary-deformation}
gives smoothness, completeness, one end with exponent $q=2>1$, and the
strictly negative boundary mean curvature. Let $\Sigma$ be its detached
minimal frontier. Its exterior is connected, complete with boundary,
and has the same end and energy. Outer area-minimization allows all
components of every enclosing competitor. Since $\Sigma$ is itself a
full cut and the metric comparison holds on its tangent three-planes,
\begin{equation}\label{four:end:area-comparison}
 |\Sigma|_{\widehat g}\ge e^{-3\eps}|\Sigma|_g
                  \ge e^{-3\eps}A_*.
\end{equation}
The numerical part of the four-dimensional Riemannian Penrose inequality
of Bray--Lee~\cite[Theorem~1.4]{BrayLee2009} now applies. Its input is a
smooth connected exterior, complete including its minimal
outer-area-minimizing compact boundary, with $R\ge0$, metric decay
$O_2(r^{-p})$, $p>1$, and scalar decay $O(r^{-s})$, $s>4$. Here the
exponents are $p=2$ and $s=4+\delta_1$; the scalar curvature is integrable
and the energy is finite. That numerical theorem permits disconnected
boundary and imposes no spin hypothesis. Its additional spin premise
concerns an equality conclusion, which is not used. Its mass normalization
is $1/(6\omega_3)$, the normalization in Equation~\eqref{four:intro:energy}.
Consequently
\[
 E_g+\frac{\Pi_N}{3\omega_3}
       (e^{-\eps N}+e^{-\eps N/2})
 \ge E_{\widehat g}
 \ge\frac12\left(\frac{|\Sigma|_{\widehat g}}{\omega_3}\right)^{2/3}
 \ge\frac12e^{-2\eps}\left(\frac{A_*}{\omega_3}\right)^{2/3}.
\]
First let $N\to\infty$ at fixed prepared data and $\eps$, then let
$\eps\downarrow0$. Only this numerical inequality is passed to the limit.
\end{proof}

\subsection{The construction and its scales}

There are four stages in the proof of the deformation theorem.
Section~\ref{four:sec:geometry} chooses two trapped-region frontiers and
their signed collars. The frontiers allow a small interval of prescribed
graph trace while the original strict dominant-energy margin pays for
the unrestricted $\tau$. Section~\ref{four:sec:system} fixes the coupled
equations and excludes first contact with the conformal floor. In
Section~\ref{four:sec:bounds}, geometric support comparisons first confine
the scaled height to the correct interval; signed boundary slopes then
make the inner flux small, before any uniform ellipticity is known.
Sections~\ref{four:sec:regularity} and~\ref{four:sec:existence} obtain
boundary regularity, construct the compact return map and complete
exhaustion and energy control.

The local scalar identities are given with their proofs and the exact
four-dimensional constants. Their common content is reusable in other
graph constructions. The degree argument, signed boundary estimates,
reflected measure estimates and four-dimensional interpolation remain
specific verifications for the boundary problem studied here.

\begin{remark}[Order of constants and limits]\label{four:rem:constants}
Fix the prepared data, the thresholds and collars, and $\eps>0$ first.
The deformation integer $N$ is then sufficiently large. A bound denoted
$\Pi_N$ is polynomial in $N$ and uniform in the outer truncation radius,
homotopy parameter and solution. Constants denoted $C(N)$ in elliptic
estimates may depend arbitrarily on $N$. Exhaust the outer radius at
fixed $N$, then pass $N\to\infty$ using only polynomial estimates, and
finally $\eps\downarrow0$. In an application to weakly decaying original
data, strictification is removed at each fixed repaired end only after
this prepared construction is finished; the end-replacement radius
tends to infinity last.
\end{remark}
